\section{Machine checks}\label{app:checks}

Apart from \Cref{prop:family2}, whose proof is the computation of
\Cref{app:cert114}, no proof in this paper depends on a computer. The finite
steps have been repeated by machine, in several independent
implementations. The programs are in the author's repository
\url{https://github.com/creelie/e84hc}, each release of which is archived on
Zenodo under the concept DOI
\href{https://doi.org/10.5281/zenodo.23234362}{10.5281/zenodo.23234362}, and
one script runs them all and compares their outputs.

\subsection*{Exact computations in three languages}
The first program exists in Python, Julia and C, all in exact integer or
rational arithmetic, and the three versions print the same
$218$ lines; the shell script compares them line by line and against a
stored copy. The lines fall into six groups.
\begin{itemize}[leftmargin=2.3em]
\item[F] Fermat varieties $X^{n}_{m}$ for $n=2$, $m\le12$ and for $n=4$,
$m\le8$: the characters of \Cref{prop:fermatchars}, the primitive Betti
number against $((m-1)^{n+2}+(-1)^{n}(m-1))/m$, the middle Hodge number, the
Hodge characters and those of pair type; for surfaces,
$h^{1,1}=(2m^{3}-6m^{2}+7m)/3$, the Picard number $\rho$, and Aoki and
Shioda's formula $\rho=3(m-1)(m-2)+1$ for $(m,6)=1$ \cite{AS83}. The values
$\rho=20,37,86$ for $m=4,5,6$ appear among them.
\item[J] Fermat curves of degree $m\le30$: every orbit of characters under
$(\ZZ/m)^{\times}$ has an odd type function, as in the proof of
\Cref{lem:fermatcurve}.
\item[S] The switch lemma (\Cref{lem:switch}) for all pairs of lists with
$p\le4$ rows and $s\le3$ columns: the column correction reaches the target
by two-row switches, each preserving the column sums, in at most
$s\lfloor p/2\rfloor$ steps.
\item[G] The signs $(-1)^{p(p-1)/2}$ of the permutations that regroup a
product of $2p$ classes of degree one between interleaved and blocked
order, for $p\le10$.
\item[W] The dimensions $3n$ and $n^{2}$ of the Prym locus and the Weil
family, and the count of \Cref{rem:count}: the pairs $(m,k)$ with
$3(m+k)+3k\ge(m+k)^{2}$, which are $k\le2$ for $m=2$, $k=0$ for $m=3$, and
none for $4\le m<60$, $k<200$; and the count of \Cref{prop:branched}: for
$2\le n\le12$ and every genus $q$, the Riemann--Hurwitz relation, the
dimension $q+2n-1$ of the family, and its unique maximum $3n$ at $q=n+1$,
$r=0$; and for \Cref{prop:abelprym}, the pairs $(a,b)$ with $0\le a\le2$,
$1\le b\le2$ and $a+b\equiv0\pmod3$, which are $(1,2)$ and $(2,1)$, the
dimensions $g$, $g-1$, $g-1$ of $V_{0}$, $V_{1}$, $V_{2}$ adding up to the
genus $3g-2$ of $C$, and the dimensions $2n^{2}$ of $E$ and $3n$ of
$H^{0}(X,K_{X}^{\otimes2})$.
\item[L] The rules of \Cref{tab:rules}: the closed sets of
\Cref{lem:closedsets}; the valuation in which \st{CM} holds and \HC{}
fails; the five minimal sets of \Cref{thm:routes}, found by enumerating all sets
of statements; and the fact that \st{L} and \st{M} form the only minimal
set avoiding \st{F2}, \st{F3$'$}, \st{IP}, \st{CM} and \st{SR}.
\end{itemize}

\subsection*{Lean}
A first Lean file is checked by the kernel of Lean~4 with the core library
only, without \texttt{sorry} and without
\texttt{native\_decide}; every theorem in it depends only on the standard
axioms. It states the rules of \Cref{tab:rules} as an inductive derivability
relation and proves: that the four sets of \Cref{lem:closedsets} and the set
$\{\st{CM},\st{Fer}\}$ are closed; \Cref{thm:routes} itself, for an arbitrary set
of statements; \Cref{cor:bypass}(i); and the derivations of \HC{} from each of
the five minimal sets. It also proves, for all sizes at once, the counting
identity behind the switch lemma, the equivalence $3n<n^{2}\iff n\ge4$,
the inequality of \Cref{rem:count} for every $m\ge4$ and $k\ge0$, and the
count of \Cref{prop:branched} for all $n$, $q$ and $r$, and the eigenspace
bookkeeping of \Cref{prop:abelprym}. By evaluation in the kernel, it
computes the Picard numbers $20$, $37$ and $86$ of the Fermat quartic,
quintic and sextic surfaces from the characters, and it checks that the sets
$O_{21}$, $O_{33}$ and $O_{39}$ of \Cref{ssec:fermatodd} are balanced,
contain no pair, generate, and have no direct or lowering move, and that the
three identities in the proof of \Cref{lem:exceptional} hold with every part
balanced.

\subsection*{Lean with Mathlib}
A second Lean project, which uses Mathlib, holds a formalization of
\Cref{thm:sextuples}, about $3600$ lines in $15$ files.
It follows the proof of \Cref{ssec:fermatfour} step by step: the counting
function, Dirichlet characters and their conductors, Fourier inversion on
$(\ZZ/e)^{\times}$, the subgroups $K_{p}$ and their orders, split functions
and \Cref{lem:grid,lem:oddbig,lem:twoD}, \Cref{prop:core},
\Cref{lem:killline,lem:killhalf}, and Step~7. Two of its theorems
state \Cref{thm:sextuples} for every $m$ prime to $6$, and a third derives the conclusion of \Cref{thm:oddfour} for
an abstract predicate ``the line $V(\alpha)$ is algebraic'' from the two
geometric inputs, algebraicity for characters that contain a pair and for
$5$-standard characters. \Cref{thm:moves}, for $m$ divisible by $3$, is not
formalized there; it is checked by the searches below for $m\le105$.

The nonvanishing of $B_{1,\chi}$ for odd primitive $\chi$ is proved there as
well, from the nonvanishing of $L(1,\bar\chi)$ in Mathlib, along the classical
route of Step~2. Expanding $\bar\chi(n)$ by the Gauss sum $\tau(\chi)$ gives
\[
  \tau(\chi)\sum_{n\ge1}\frac{\bar\chi(n)}{n}
  =-\sum_{s=1}^{e-1}\chi(s)\log\bigl(1-e^{2\pi is/e}\bigr)
  =-\frac{\pi i}{e}\sum_{s=1}^{e-1}\chi(s)\,s,
\]
where the logarithmic series converges on the unit circle away from $1$ by
Dirichlet's test and Abel's limit theorem, the real parts cancel because
$\chi$ is odd, and $\arg(1-e^{i\theta})=\theta/2-\pi/2$ for
$0<\theta<2\pi$. A summation by parts, uniform in $\sigma\in[1,2]$, shows that
the value at $s=1$ of the analytic continuation of $L(s,\bar\chi)$ is the sum
of the conditionally convergent series on the left. The geometric inputs and
the description of the Hodge classes by characters remain hypotheses of the
formal statements, since Mathlib has no Hodge theory. With them, the main
theorems, among them the classification and the conclusion of \Cref{thm:oddfour}
with $B_{1,\chi}\ne0$ discharged, depend only on the standard axioms, and no
file uses \texttt{sorry}.

\subsection*{Fermat fourfolds by exhaustive search}
A second program exists in Python, Julia and C. For
every $m$ prime to $6$ up to a bound it lists the balanced multisets of four
and of six nonzero residues modulo $m$ and checks \Cref{thm:sextuples}: every
quadruple contains a pair, and every sextuple without a pair is $5$-standard.
The three versions print the same lines for $m\le55$, and the C version is run
up to $m=125$ and compared with a stored copy of its output; for $m=125$, for
instance, it finds $41688$ balanced sextuples, of which $24$ contain no pair,
all of them $5$-standard.

A third program, in the same three languages, treats the odd
$m$ divisible by $3$. It lists the balanced sextuples, and for each one that
contains no pair and generates $\ZZ/m$ it finds a direct or a lowering move
(\Cref{def:move}), or checks that $tA$ is one of $O_{21}$, $O_{33}$, $O_{39}$
for a unit $t$, as \Cref{thm:moves} asserts. It also checks the identities of
\Cref{lem:exceptional} and that the three exceptional sets have no move. The
three versions print the same lines for $m\le63$, and the C version is run up
to $m=105$ and compared with a stored copy of its output; for $m=105$ it
finds $27690$ balanced sextuples, of which $78$ contain no pair and $54$ of
these generate $\ZZ/105$, each with a direct move. These searches check the
theorems only for the degrees listed; the proofs cover every $m$.

\subsection*{Macaulay2}
Two scripts use Macaulay2. The first computes the primitive Hodge numbers of
the Fermat surfaces and fourfolds from the Jacobian ring,
$h^{n-q,q}_{\mathrm{prim}}=\dim R_{(q+1)m-(n+2)}$ with
$R=\QQ[x_{0},\dots,x_{n+1}]/(x_{0}^{m-1},\dots,x_{n+1}^{m-1})$
\cite{Gri69}, and checks that they agree with the character counts of
item~F. The second takes random smooth plane curves of degree $4\le d\le8$
modulo the prime $32003$ and checks Max Noether's theorem in the form used in
\Cref{ssec:whereitstops}: the multiplication map
$\operatorname{Sym}^{2}H^{0}(K)\to H^{0}(2K)$ has rank $3g-3$, so
$\dim I_{2}=(g-2)(g-3)/2$. A smooth curve of full rank modulo $p$ is smooth
and of full rank over $\QQ$, so these are checks over $\QQ$.

\subsection*{Degree \texorpdfstring{$114$}{114}}
A program in Python, Julia and C checks the finite statements of
\Cref{sec:fermat114}: \Cref{lem:twochars}, by listing $\langle t\alpha_{1}
\rangle$ and $\langle t\alpha_{2}\rangle$ for the $36$ units $t$; that
$\nu_{19}$ is $1$ on $\beta$ and vanishes on every pair and on the $98$
standard multisets of level $114$; and the weights and degrees of the two
families. The Python and Julia versions also check \Cref{prop:family1} in
exact rational arithmetic at three values of $\lambda$, and a computer
algebra system confirms the closed formula for $I_{20}$. The lattice
computation described in \Cref{app:cert114} is done with exact integer
matrices. In the Lean project with Mathlib, a further file proves
\Cref{lem:twochars}(a),(b), the vanishing of $\nu_{19}$ on the pairs and on
the standard multisets, the identities in the proof of \Cref{prop:family1}
(that $\sum_{k}u_{k}=0$, the two residues, the closed formula and its value
at $\lambda=2$) and the weight and degree conditions of both families; its
theorems depend only on the standard axioms. The certificate of
\Cref{app:cert114} is run at $1000$ and at $600$ bits.

\subsection*{References}
A script looks up every entry of
the bibliography online, through its DOI where it has one and through
Crossref, DataCite, arXiv or the publisher's page otherwise, and compares
title, year, volume and pages; its report is stored next to it.

Except for the certificate of \Cref{app:cert114}, none of these checks
touches the geometry the paper relies on, and none of them bears on
\Cref{q:schoen}.
